\problema{Number of Islands}{200}{src/01-grafos/200-number-of-islands.cpp}{M}

Nos dan una cuadrícula de caracteres donde \texttt{'1'} es tierra y
\texttt{'0'} es agua. Una \emph{isla} es un grupo de tierras conectadas en
horizontal o vertical. Queremos contar cuántas islas hay.

\paragraph{Intuición.} Imagina que sobrevuelas el mapa. Cada vez que ves un
pedazo de tierra que aún no has pisado, plantas una bandera (una isla nueva)
y luego caminas por toda la tierra conectada hundiéndola en el mapa, para no
volver a contarla. Repites hasta que no quede tierra nueva.

\begin{center}
\ttfamily
\begin{tabular}{ccccc}
1 & 1 & 0 & 0 & 0\\
1 & 1 & 0 & 0 & 0\\
0 & 0 & 1 & 0 & 0\\
0 & 0 & 0 & 1 & 1\\
\end{tabular}
\end{center}

\noindent En esta cuadrícula hay \textbf{3} islas: el bloque $2\times2$ de
arriba a la izquierda, la celda central, y el par de abajo a la derecha.

\paragraph{Solución.} Recorremos la cuadrícula; al encontrar tierra nueva
sumamos una isla y usamos DFS (\texttt{hundir}) para marcar como agua toda su
tierra conectada. El código completo (abre el archivo con clic en el título):

\lstinputlisting{src/01-grafos/200-number-of-islands.cpp}

\complejidad{$O(m\cdot n)$}{$O(m\cdot n)$}

\paragraph{Notas de entrevista (Amazon).} Es el clásico de Amazon. Menciona
que modificar la cuadrícula in situ evita el arreglo \texttt{visited}; si no
puedes mutar la entrada, usa un conjunto de visitados o BFS con cola. Casos
borde: cuadrícula vacía, toda agua, toda tierra (una isla).
